%------------------------------------------------------------------------------%
\begin{question}
  Calcule
  \[
    \int_{0}^3 \frac{1}{x-1} \, dx
  \]
  se possível
\end{question}

%------------------------------------------------------------------------------%
\begin{resolution}{Solução}
  A função \(f(x) = \frac{1}{x-1}\) é descontínua em $x=1\in[0, 3]$

  \pause
  \bigskip
  Precisamos dividir o intervalo e duas partes \([0, 1]\) e \([1, 3]\)
  \pause
  \[
          \int_0^3 \frac{1}{x-1} \,dx
    \;=\; \int_0^1 \frac{1}{x-1} \,dx
    \;+\; \int_1^3 \frac{1}{x-1} \,dx
  \]

  \pause
  Agora a descontinuidade está nos extremos dos intervalos
  e podemos usar integrais impróprias
\end{resolution}

%------------------------------------------------------------------------------%
\begin{resolution}{Solução}
\begin{align*}
  \onslide<+->{\int_{0}^1 \frac{1}{x-1} dx}
  \onslide<+->{& = \lim_{t \to 1^{-}}\int_{0}^t \frac{1}{x-1} dx  \\}
  \onslide<+->{& = \lim_{t \to 1^{-}} \ln\abs{x-1} \evalat{0}{t}  \\}
  \onslide<+->{& = \lim_{t \to 1^{-}}\big(\ln\abs{t-1} - \ln\abs{-1} \big) \\}
  \onslide<+->{& = \lim_{t \to 1^{-}} \ln(1-t) }
  \onslide<+->{& & 1-t \to 0^+ \text{ quando } t \to 1^{-} \\}
  \onslide<+->{& =-\infty}
\end{align*}
\onslide<+->{Portanto a integral diverge e não precisamos calcular a segunda parte}
\end{resolution}

%------------------------------------------------------------------------------%